% Conclusions -- budget 210 words (body text).
% Sources: Sections 1, 6 and 7 of this paper; docs/results_summary.md (macro-F1 0.90-0.94, ECE 0.07-0.10, outage counts
% closer in 30 of 30 datasets, Tetouan A1); limitations of Section 7.3; deliberate manipulation as future work at the
% author's request (answer 1 on Section 2; hassine2025security verified, notes in references.bib).
\section{Conclusions}\label{sec:conclusions}

This paper typed the anomalies of a multi-city feeder record as acquisition artefacts or grid events, with calibrated probabilities and without labelled data. Rules read the level of each anomaly and the traces a recording chain leaves, a constrained label model turns their sparse votes into stable probabilistic labels, and a residual detector adds anomalies the rules cannot type. By our typing, most of what looks like an outage in the Moroccan telemetry is an acquisition failure; straight-line interpolation is the most common fill, and restoration shape carries almost no evidence. On injected events the labels reach macro-F1 0.90--0.94 at an expected calibration error of 0.07--0.10, typed counting is closer to the true outage time than naive counting in every dataset, and the method transfers to another city without re-tuning, except for zero-fill gaps, whose evidence is cross-city. The released catalogue lets later studies impute the artefacts and keep the grid events.

Four extensions follow. Utility outage records would validate the labels on real events. Untyped detector candidates, partial drops, surges and short disturbances, call for new types. An episode-level label model would use context that item-level labels ignore. Deliberately manipulated readings, a concern of smart-grid security \cite{hassine2025security}, would add a third cause beside artefacts and grid events.
